\chapter{Amulet practices across different social, temporal and regional contexts}

The chapter presents the comparative analysis and interpretation of the different objectscapes established in the previous chapters. The amulet motifs were categorized into interpretative subcategories to support the analysis. Additionally, since many elite and non-elite tombs were dated over a broad range of phases (e.g. 25th Dynasty-Middle Napatan, 25th Dynasty-Early Napatan, etc.). In order to improve direct comparisons, the data was divided into earlier and later phases. Material composition will also be analyzed. These were divided into local and imported materials. The percentage, however, does not total 100\%, as some materials were considered as both local and imported when there were sources in both Sudan and Egypt. Furthermore, it is difficult to ascertain whether the materials or finished amulets were imported, as available evidence comes only from the Sanam Treasury and temple workshop.

The data will be analyzed across three main elements: social, chronological and regional. The center of the analysis is the social groups, which will then be evaluated within different chronological phases and locations. The chronological division follows the same one established in previous chapters. The regional division is as follows, from north to south of the Sudanese territory: Lower Nubia (Mirgissa, Qustul and Debeira 176), North Upper Nubia (Missiminia), 4th Cataract (El-Kurru, Nuri and Sanam) and Meroe region (Meroe West and South).

Apart from the internal comparisons between the sites from the database, the amulets will also be compared with other local and foreign examples. In the local perspective, recently published contemporary sites from Sudan, namely Kawa and Sedeinga, can contribute to establish both objectscapes within Nubia as well as internal diversity. Meanwhile, contemporary amulets from ancient Egypt, Near East and Mediterranean can be used to compare practices between cultures and assess the extent of influence and spread of Egyptian cultural markers.

\section{Chronological variations}

This section focus on the chronological variations of amulets across royal, elite and non-elite assemblages. The main objective is to interpret these variations based on the historical context of each phase and the internal social heterogeneity within Nubia. As a result, we can establish specific objectscapes of each different context in relation to amulets.

\subsection{Amulet distribution}

First, it is important to understand the distribution of amulets, especially their frequency in burials. As shown in Figure \ref{fig:phase-distro}, amulets were not included in the majority of burials, regardless of social group, with elite and non-elite contexts displaying the lowest distribution. Only in royal sites of the 25th Dynasty and Early Napatan phases more than half tombs contained at least one amulet―70\% and 52\%, respectively. In fact, royal contexts display the biggest percentage of tombs with amulets, with the exception of those from the pre-25th Dynasty (8\%).
In contrast, less than 12\% of elite tombs from all phases contained any amulets, with the highest percentage stemming from Late Napatan sites. For the non-elite, distribution follows the royal pattern, with more amulets present in tombs dated to the 25th Dynasty and Early Napatan phases (24\% and 22\%). Fewer tombs from the remaining phases contained amulets, only 5\% from pre-25th Dynasty and 6\% from Middle Napatan.

\begin{figure}
	\centering
	\includegraphics[width=1\textwidth]{phase-distro}
	\caption{Distribution (in percentage) of tombs containing amulets by social group and chronological phase.}
	\label{fig:phase-distro}
\end{figure}

This phenomenon could have been affected by the level of disturbance of the tombs, as the vast majority were found already looted, especially the royal pyramids. Therefore, it is interesting to look into the intact tombs, which were located only at the Meroe cemeteries, as noted in the reports by Reisner/Dunham. In the case of Sanam, Griffith mentions that the “plundering that took place in antiquity was extensive amongst all except the poorest graves” \parencite[p.~75]{griffithOxfordExcavationsNubia1923}, without presenting numbers or mentioning specific tombs. Therefore, it is impossible to assess with any level of certainty how many tombs were undisturbed.

Taking into account the intact tombs from Meroe, it appears that non-elite graves were not a priority for plunderers (Figure \ref{fig:intact-meroe}). This can be attributed to their lack of superstructures, which would make than less visible, and possibly less desirable. There were nine intact burials from the elite, with one dated to the 25th Dynasty and with amulets. In contrast, out of eight total intact tombs from the Early Napatan period, there were only two examples where amulets were found (25\%). In non-elite intact tombs, amulets were less frequent, although the number of intact tombs was higher. Less than 20\% of 25th Dynasty graves contained amulets, while for Early Napatan the percentage was 15\% and Middle Napatan, 19\%. Considering these figures, there is reason to believe that amulets were not an integral part of the funerary equipment, especially for the elite and non-elite. It is possible that their inclusion was based on personal choice.

\begin{figure}[H]
	\centering
	\includegraphics[width=1\textwidth]{intact-meroe}
	\caption{Comparison of total number of intact tombs and those containing amulets in Meroe cemeteries by social group and chronological phase.}
	\label{fig:intact-meroe}
\end{figure}

\subsection{Amulet types}

\begin{figure}
	\centering
	\includegraphics[width=1\textwidth]{phase-types}
	\caption{Frequency (in percentage) of amulet categories by social group and chronological phase.}
	\label{fig:phase-types}
\end{figure}

Regarding the distribution of amulet types, certain patterns and differences can be observed across social groups (Figure \ref{fig:phase-types}). In royal contexts, symbols are the most popular category of amulets in all phases. They were mostly followed by deities―except during the Middle Napatan phase, when animals came in second. Some categories are totally absent in a given phase. For instance, amulets of animals, objects and nature were not found in pre-25th Dynasty and Late Napatan burials, whereas object amulets were absent from Middle Napatan graves. The most diverse assemblages come from 25th Dynasty and Early Napatan contexts, encompassing almost all seven categories, with the exception of two types: miscellaneous motifs were not included in Early Napatan burials, while amulets of humans were absent from Late Napatan graves. In contrast, pre-25th Dynasty and Middle Napatan assemblages included three amulet types, while Late Napatan was the least diverse, with only amulets of symbols and deities.

Similarly, symbol amulets were the most common in elite tombs throughout the Napatan period, with deities and animals appearing as the second and third most frequent categories from the 25th Dynasty to Middle Napatan phases. Meanwhile, Late Napatan assemblages were the least diverse, containing only symbols and animals. Six of the seven amulet types were present in 25th Dynasty and Middle Napatan assemblages. The earlier phase did not contain any amulets of objects, while human figures were absent from the latter. Early Napatan amulets were the most diverse, with all categories found.

Non-elite assemblages differ significantly from the other social groups, especially during the pre-25th Dynasty and Middle Napatan phases, when deities were the most common amulet type. Additionally, the second most common type during the Middle Napatan phase were human figures and not symbols, as is the case in the pre-25th Dynasty. 25th Dynasty and Early Napatan amulets follow the same pattern as those from contemporary royal and elite, when the three most popular types being symbols, deities and animals, in this order. Both phases were also the most diverse, including all seven amulet categories. In contrast, Middle Napatan burials included six amulet types, while those from the pre-25th Dynasty had only three; symbols, deities and animals. Compared to elite burials, there is an increase in the number of human figure amulets in non-elite contexts starting in the 25th Dynasty period and peaked in the Middle Napatan, being the second most common type.

Important insights can be derived from this data. Firstly, the prevalence of symbol amulets is mainly observed in elite and royal burials. Additionally, there appears to be an increase in deity amulets within non-elite graves. Furthermore, human figures were found only in elite and non-elite context and, in these contexts, there is a higher frequency of animal motifs. Finally, assemblages from 25th Dynasty and Early Napatan across all groups display a greater variety of types. In subsequent sections, we will explore these observations in closer detail with the analysis of subcategories of each amulet type in order to further interpret their use in certain contexts.

\subsection{Amulet motifs}

\begin{figure}
	\centering
	\includegraphics[width=1\textwidth]{phase-deities}
	\caption{Proportion (in percentage) of amulets of Egyptian deities and local adaptations by social group and chronological phase, relative to the total number of amulets.}
	\label{fig:phase-deities}
\end{figure}

As previously mentioned, in order to better analyze the distribution of certain motifs, they were divided into interpretative subcategories based on their broader types. This choice was made with the aim to explore consumption patterns and differences between social groups, as well as examining the impact of changing relations with Egypt during the Napatan period.


Deities amulets were divided into two categories: those from the Egyptian pantheon and local adaptations and/or deities. Egyptian deities as considered those that commonly appear in the religious sphere, while local adaptations/deities are depictions of deities that differ from the standard Egyptian design―although oftentimes they were paired with Egyptian symbols. Among those are the group of winged goddesses and gods, a ram-headed dwarf, unidentified deities―such as those with snake, hawk, double hawk, crocodile and ram heads―, and adaptations of the Isis nursing Horus motifs―for instance, the lion-headed goddess nursing Horus and Isis nursing a queen.

Overall, Egyptian deities were the most widespread across all phases and groups, but local adaptation/deities presents varied significantly (Figure \ref{fig:phase-deities}). In royal contexts, during the pre-25th Dynasty, only a single unidentified hawk-headed dwarf back-to-back with a Pataikos is a local motif; whereas the 25th Dynasty amulets included more local motifs, comprising 8\% of the total amulets. No local motifs were found in the subsequent phases in royal burials. In contrast, in elite graves, local adaptation/deities found were found during the Early and Middle Napatan periods―1\% and 6\% of all amulets, respectively. These included two unknown deities, one with a hawk head and another with a snake head.


For the non-elite, the occurrence of these amulets was even rarer, less than 1\% of all amulets in both the 25th Dynasty and Early Napatan phases. The motifs encompass the lion-headed goddess nursing Horus, deities with the heads of a hawk, lion, snake and crocodile, a lion-headed goddess, two winged lion-headed goddess and a winged anthropomorphic goddess. Although the absolute quantity of these amulets is higher than in other social groups, the total amount of non-elite amulets is so high that their percentage is not considerable.


The data indicates that local adaptation/deities were proportionally more common in royal contexts of the 25th Dynasty, and especially reserved to queens and minor wives. However, certain individuals from the elite and non-elite could also acquire these kind of amulets and transform them into grave goods.

\begin{figure}
	\centering
	\includegraphics[width=1\textwidth]{phase-udjat}
	\caption{Proportion (in percentage) of \textit{udjat} and quadruple \textit{udjat} amulets by social group and chronological phase, relative to the total number of amulets.}
	\label{fig:phase-udjat}
\end{figure}

Amulets of symbols were divided into two graphs, as those depicting \textit{udjat} were one of the most prevalent motifs (Figure \ref{fig:phase-udjat}). Quadruple udjat amulets were also included in the same chart as the singular \textit{udjat}. Proportionally, \textit{udjat} amulets were mainly popular during the 25th Dynasty in all social groups. Amulets from the sole pre-25th Dynasty royal burial were all \textit{udjat}, with a sole exception.


From the Early Napatan royal contexts, the quantity of \textit{udjat} amulets decreases, reaching the lowest proportion in the Middle Napatan phase. However, it spikes again in Late Napatan tombs, becoming proportionally higher than on the 25th Dynasty. Elite \textit{udjat} amulets display a clear decrease in use after the peak during the  25th Dynasty. On the other hand, non-elite burials displayed the least quantity of \textit{udjat} during the pre-25th Dynasty and Middle Napatan phases, with the highest amount present in 25th Dynasty and Early Napatan graves. These contexts were also the only ones that included quadruple \textit{udjat} amulets. Although, they were not found in great quantities―2\% and 5\%, respectively―, it is interesting that only the non-elite consumed these amulets.


The prevalence of \textit{udjat} amulets follows a similar pattern attested in other contexts, such as the Nile Delta, Upper Egypt, Achaemenid Persia, the Levant, and Mediterranean; spanning the Third Intermediate and Late Periods \parencites{vergaConsiderazioniMargineSignificato1981, herrmannÄgyptischeAmuletteAus1994, herrmannÄgyptischeAmuletteAus2002, herrmannÄgyptischeAmuletteAus2016}[p.~374-376]{astonBurialAssemblagesDynasty2009}{lópez-grandeAmuletosIconografíaEgipcia2014}{arroyoPapelDosAmuletos2020, arroyoUsoAmuletosComo2021}{qahériObjectsÉgyptiensÉgyptianisants2020}{muñozpérezEgyptianAmuletsWestern2023}{furlanAmuletExchangeNile2024}. Accordingly, the evidence from Napatan Nubia is not exceptional. However, the setting differs significantly from other contexts. After nearly 500 years of Egyptian occupation, the degree of entanglement between both cultures likely resulted in Nubians not considering Egyptian and Egyptian-style objects as foreign by the Nubians but rather as part of their cultural landscapes and everyday life. Therefore, the use of \textit{udjat} amulets CONTINUE

\begin{figure}
	\centering
	\includegraphics[width=1\textwidth]{phase-symbols}
	\caption{Proportion (in percentage) of symbol amulets, divided in subcategories, by social group and chronological phase, relative to the total number of amulets.}
	\label{fig:phase-symbols}
\end{figure}

Regarding the additional symbols amulets, they were separated into subcategories relative to their motifs (Figure \ref{fig:phase-symbols}). These include winged motifs, such as winged griffins, scarabs, snakes and uraeus; hieroglyphic signs representing Egyptian words (ankh, \textit{nefer}, \textit{sa}, \textit{sema}, \textit{nt}, \textit{nb}, basket and raised arms); royal symbols encompassing the crook, was scepter, whip, \textit{uraeus}, double \textit{uraeus}, \textit{hmhm} crown and double feather; symbols with animals heads, such as \textit{uraei} with lion heads and ram-headed sphinxes; symbols related to the sun include the \textit{akhet} (sun horizon), sun boat, sun disc, and horned sun disc; moon-related symbols (horned lunar disc and lunar crescent); common symbols which are very recurrent in amulets across Egypt and Nubia, such as the \textit{menat}, heart, \textit{djed}, and Isis knot; and finally, single and double \textit{ba} birds and sphinxes.


Royal burials contained, proportionally, the highest amount of these symbol amulets. As previously mentioned, the sole pre-25th Dynasty burial with amulets contained only \textit{udjat}. Conversely, 25th Dynasty tombs included the most diverse assemblage, with seven out of nine subcategories represented―sun and moon-related motifs were absent. The majority of amulets represented winged motifs and common symbols. On the other hand, Early Napatan symbol amulets depicted only common and royal symbols, while those from Middle Napatan included mostly common symbols, but also hieroglyphic signs,, ba birds, and royal and sun-related motifs. In all these phases, common symbols were the most popular motif, with the \textit{menat} being particularly prominent, as mentioned in Chapter 4. However, this pattern changes in Late Napata, when royal symbols were slightly more common than the other concurrent subcategories―winged and common motifs, and ba birds.


Elite contexts follow the overall pattern, with common motifs the most prominent in all phases, although they represented less than 40\% of the total amulets in this group. In fact, during the 25th Dynasty and Late Napatan, they were the only symbol amulets, and making up 40\% of the total amulets in the later phase. Unlike royal contexts, only only one \textit{menat} was found in elite burials. Most of the common motifs were composed of \textit{uraei}, hearts and \textit{djed}.


Apart from \textit{udjat} and quadruple \textit{udjat}, symbol amulets consumption patterns of the non-elite are completely different from the other social groups. Firstly, they were only found in burials dated to the 25th Dynasty and Early Napatan phases. Secondly, they represented less than 10\% of the total amulets in each phase, and finally, royal motifs were the most frequent during the 25th Dynasty, while in the Early Napatan period, royal motifs and hieroglyphic signs were slightly more recurring the other motifs. Other subcategories were present but in significantly lower quantities, some with a unique example.

\begin{figure}
	\centering
	\includegraphics[width=1\textwidth]{phase-animals}
	\caption{Proportion (in percentage) of symbol amulets, divided in subcategories, by social group and chronological phase, relative to the total number of amulets.}
	\label{fig:phase-animals}
\end{figure}

Animal amulets were divided into six subcategories (Figure \ref{fig:phase-animals}):


- Animals associated with common Egyptian deities. In the corpus, these include cats (Bastet), jackals (Anubis), ibises and apes (Thoth), crocodiles (Sobek), hippopotamus (Taweret), lions (Maahes, Sekhmet and Tefnut), rams (Amun), snakes (Nehebkau), and falcons (Horus). Although, it is also important to consider that many of these animals could be found in the Nubian nature and, therefore, were not necessarily associated with an Egyptian deity, but stem from observation.


- Common animals, which are not associated with popular deities from the Egyptian pantheon, such as unidentified birds, ducks, fishes, flies, frogs, giraffes, geese, hares, hedgehogs, lizards, monkeys, pigs, scorpions, sows and turtles.


- Double animals include double frog and double ram motifs.


- Quadruple animals encompasses four apes and a four-headed ram amulets.


- Hybrid animals include amulets that combine elements from different animals, such as hawk-headed crocodiles, ram-headed scarabs, and lion-headed flies.
- Parts of animals consists of feathers and vulture wings motifs.


With a few exceptions, animals associated with Egyptian deities were the majority in all social groups. In royal contexts, only three subcategories were present in three phases―no animal amulets were found during the pre-25th Dynasty and Late Napatan phases. For 25th Dynasty and Early Napatan amulets, only motifs of animals associated with Egyptian deities and common animals were present, with the first representing the majority. Although in Middle Napatan, animals related to Egyptian deities were still the most popular, no common animals were present, being replaced by quadruple animals, in this case, a four-headed ram.

Although animal amulets were more prevalent in elite contexts, they were limited to those linked to Egyptian deities and common animals, mirroring the pattern of early royal assemblages. However, in 25th Dynasty elite burials, amulets of common animals were more prevalent than those linked to deities. Similar to royal assemblages, during the Early and Middle Napatan periods, amulets associated with deities outnumbered those of common animals. Conversely, Late Napatan contexts contained only animals linked to deities.

There is a marked shift in the pattern for non-elite assemblages, after the pre-25th Dynasty period, when only animals of Egyptian deities were found. Especially during the 25th Dynasty and Early Napatan periods, not only do common animals exceed deity-associated ones, but more subcategories are present. These include hybrid animals, and double and quadruple animals. Additionally, Early Napatan assemblages also contained amulets depicting parts of animals. Common animal amulets were also more popular than those linked to deities during the Middle Napatan phase.

\begin{figure}
	\centering
	\includegraphics[width=1\textwidth]{phase-nature}
	\caption{Proportion (in percentage) of nature amulets, divided in subcategories, by social group and chronological phase, relative to the total number of amulets.}
	\label{fig:phase-nature}
\end{figure}


\begin{figure}
	\centering
	\includegraphics[width=1\textwidth]{phase-human}
	\caption{Proportion (in percentage) of human amulets, divided in subcategories, by social group and chronological phase, relative to the total number of amulets.}
	\label{fig:phase-human}
\end{figure}


\begin{figure}
	\centering
	\includegraphics[width=1\textwidth]{phase-objs-misc}
	\caption{Proportion (in percentage) of object and miscellaneous amulets, divided in subcategories, by social group and chronological phase, relative to the total number of amulets.}
	\label{fig:phase-objs-misc}
\end{figure}


\begin{figure}
	\centering
	\includegraphics[width=1\textwidth]{phase-2deities}
	\caption{Distribution of deity amulets with two or more motifs by social group and chronological phase.}
	\label{fig:phase-2deities}
\end{figure}

\subsection{Amulet materials}

\begin{figure}
	\centering
	\includegraphics[width=1\textwidth]{phase-faience}
	\caption{Percentage of faience amulets by social group and chronological phase.}
	\label{fig:phase-faience}
\end{figure}

\begin{figure}
	\centering
	\includegraphics[width=1\textwidth]{phase-materials}
	\caption{Percentage of local and imported materials by social group and chronological phase.}
	\label{fig:phase-materials}
\end{figure}


\begin{figure}
	\centering
	\includegraphics[width=1\textwidth]{phase-materials-selec}
	\caption{Percentage of selected materials―gold, silver, lapis lazuli, ivory, electrum and obsidian―by social group and chronological phase.}
	\label{fig:phase-materials-selec}
\end{figure}

\section{Regional variations}

In this section, amulets will be analyzed based on the geographical location of the sites, in order to build regional objectscapes. Elite and non-elite burials were distributed across different locations in the Nubian territory, while royal tombs were only found in the 4th Cataract region, where the center of the Napatan period was located.

\subsection{Amulet distribution}

\begin{figure}
	\centering
	\includegraphics[width=1\textwidth]{region-distro}
	\caption{Distribution (in percentage) of tombs containing amulets by social group and region.}
	\label{fig:region-distro}
\end{figure}

\subsection{Amulet types}

\begin{figure}
	\centering
	\includegraphics[width=1\textwidth]{region-types}
	\caption{Frequency (in percentage) of amulet categories by social group and region.}
	\label{fig:region-types}
\end{figure}


\subsection{Amulet motifs}

\begin{figure}
	\centering
	\includegraphics[width=1\textwidth]{region-deities}
	\caption{Proportion (in percentage) of amulets of Egyptian deities and local adaptations by social group and region, relative to the total number of amulets.}
	\label{fig:region-deities}
\end{figure}


\begin{figure}
	\centering
	\includegraphics[width=1\textwidth]{region-udjat}
	\caption{Proportion (in percentage) of udjat and quadruple udjat amulets by social group and region, relative to the total number of amulets.}
	\label{fig:region-udjat}
\end{figure}


\begin{figure}
	\centering
	\includegraphics[width=1\textwidth]{region-symbols}
	\caption{Proportion (in percentage) of symbol amulets, divided in subcategories, by social group and region, relative to the total number of amulets.}
	\label{fig:region-symbols}
\end{figure}


\begin{figure}
	\centering
	\includegraphics[width=1\textwidth]{region-animals}
	\caption{Proportion (in percentage) of animal amulets, divided in subcategories, by social group and region, relative to the total number of amulets.}
	\label{fig:region-animals}
\end{figure}


\begin{figure}
	\centering
	\includegraphics[width=1\textwidth]{region-nature-human}
	\caption{Proportion (in percentage) of nature and human amulets, divided in subcategories, by social group and region, relative to the total number of amulets.}
	\label{fig:region-nature-human}
\end{figure}


\begin{figure}
	\centering
	\includegraphics[width=1\textwidth]{region-objs-misc}
	\caption{Proportion (in percentage) of object and miscellaneous amulets, divided in subcategories, by social group and region, relative to the total number of amulets.}
	\label{fig:region-objs-misc}
\end{figure}


\begin{figure}
	\centering
	\includegraphics[width=1\textwidth]{region-2deities}
	\caption{Distribution of deity amulets with two or more motifs by social group and region.}
	\label{fig:region-2deities}
\end{figure}


\subsection{Amulet materials}


\begin{figure}
	\centering
	\includegraphics[width=1\textwidth]{region-faience}
	\caption{Percentage of faience amulets by social group and region.}
	\label{fig:region-faience}
\end{figure}


\begin{figure}
	\centering
	\includegraphics[width=1\textwidth]{region-materials}
	\caption{Percentage of local and imported materials by social group and region.}
	\label{fig:region-materials}
\end{figure}


\begin{figure}
	\centering
	\includegraphics[width=1\textwidth]{region-materials-selec}
	\caption{Percentage of selected materials―gold, silver, lapis lazuli, ivory, electrum and obsidian―by social group and region.}
	\label{fig:region-materials-selec}
\end{figure}


\section{Conclusion}